\subsection{Soundness of the Primitive Timed Translations}
\label{subsec:timed-soundness}

For each primitive hatted timed constructor $F$, the local soundness lemma has
the form
\begin{equation}
  \rho,i\satL \T_\pi(F)
  \quad\Longrightarrow\quad
  \ewbase(\rho),i\satM F,
\label{eq:local-soundness}
\end{equation}
assuming clock consistency and the induction hypotheses for the immediate
subformulas.  No canonical reset-policy assumption appears in
\eqref{eq:local-soundness}.

\runinhead{Upper-bounded hatted Until.}
The translated LTL Until starts at $i+1$ and provides an endpoint $j>i$.
At position $i$ itself the translation does not constrain the clock, but the
recurrence~\eqref{eq:clock-update} together with nonnegativity gives
$\CVal_\rho(i+1,x)\ge\Delta_i$ whether or not $x$ is reset at $i$.  Along the
intermediate positions $(i,j)$ the shared clock is unchanged, so
Lemma~\ref{lem:clock-accumulation} applied on $[i+1,j)$ yields
$\CVal_\rho(j,x)\ge\Delta_i+\elapsed_w(i+1,j)=\elapsed_w(i,j)$, and the
endpoint guard $x\le d$ gives
\[
  \elapsed_w(i,j)\le d.
\]
The induction hypotheses turn the translated prefix $\overline{P}$ into $p$ and the
endpoint $\overline{Q}$ into $q$, matching the semantics of
$p\hUntil_{\le d}q$.  For $\hUntil_{<d}$, the strict guard $x<d$ gives
$\elapsed_w(i,j)<d$ in the same way (\texttt{sound\_\allowbreak{}UhatLt}).

\runinhead{Lower-bounded hatted Until.}
The translation resets $x$ at $i$ and follows $\Gamma$.  In the finite branch,
a future state satisfies $x\ge d$ together with $\overline{P}\Until \overline{Q}$.  By
\eqref{eq:lower-clock}, the endpoint is at least $d$ time units after $i$, and
the nested Until supplies the required $p$-prefix and $q$-endpoint.  In the
second branch, $\Always \overline{P}\land\Always\Eventually \overline{Q}$ holds.  Then $p$ holds
forever and $q$ recurs infinitely often; time divergence guarantees a $q$
occurrence at physical distance at least $d$.  For $\hUntil_{>d}$, the strict
threshold $x>d$ gives a strict lower bound in the finite branch, and the
fairness branch uses the strict argument of Section~\ref{subsec:main-strict-bounds}
(\texttt{sound\_\allowbreak{}UhatGt}).

\runinhead{Upper-bounded hatted Release.}
The translation resets the shared clock and requires
\[
  \overline{Q}\WeakUntil\bigl((x>d)\lor(\overline{P}\land \overline{Q})\bigr).
\]
Before the weak-until endpoint, $q$ holds at every strictly future position.
If the endpoint is $\overline{P}\land \overline{Q}$, $p$ discharges the Release obligation; if it is
$x>d$, the upper-bounded window has ended.  For $\hRelease_{<d}$, the window
ends at $x\ge d$ (\texttt{sound\_\allowbreak{}RhatLt}).

\runinhead{Lower-bounded hatted Release.}
The translation requires
\[
  ((x<d)\land\Unch(x))
  \WeakUntil
  \bigl((\overline{P}\Release \overline{Q})\lor((x<d)\land \overline{P})\bigr)
\]
from $i+1$.  While the left side holds, the clock is preserved and remains
below $d$, so those positions are still before the lower threshold.  At the
endpoint either an ordinary untimed Release takes over or $p$ discharges the
obligation before the threshold is reached.  For $\hRelease_{>d}$, the
pre-threshold test is $x\le d$ (\texttt{sound\_\allowbreak{}RhatGt}).

There are no separate ordinary timed soundness lemmas.  After unfolding
Eqs.~\eqref{eq:derive-ule}--\eqref{eq:derive-rge}, their
correctness follows from the Boolean cases of the structural induction and the
eight hatted lemmas above.

\subsection{Completeness under the Canonical Reset Policy}
\label{subsec:timed-completeness}

The converse local result is proved for the canonical extension:
\[
  w,i\satM F
  \quad\Longrightarrow\quad
  \rho^\star,i\satL\T_\pi(F).
\]
The difficult point is that all occurrences of the same primitive timed formula
share one clock and each occurrence may itself be activated repeatedly.  Thus
several semantic obligations can compete for the same clock trace.  The
canonical reset policy and the residual predicates are designed to resolve
this interaction.

\runinhead{Upper-bounded hatted Until.}
Choose the first $q$-position among the MTL witnesses.  The proof maintains,
for every position $n$ with $i<n\le j$ up to that endpoint $j$, the invariant
\[
  \CVal_{\rho^\star}(n,x)+\elapsed_w(n,j)\le d.
\]
It need not hold at the activation position $i$ itself (the policy may reset
there because the old value is useless); it is established at $i+1$ in both
the reset and the preservation case.
If the policy preserves the clock, Section~\ref{subsec:canonical-main}
guarantees an older reference whose residual $\mathsf{URes}$ is still
satisfiable strictly after the current position; if it resets, the new value is
consistent with the current activation.  Hence every intermediate position satisfies the clock
guard required by the translated Until and the first $q$-position supplies its
endpoint.  This is the content of the mechanized lemma
\texttt{complete\_\allowbreak{}UhatLe}.  The strict case maintains the strict
invariant $\CVal_{\rho^\star}(n,x)+\elapsed_w(n,j)<d$ with the residual
$\mathsf{URes}^{<}$ (\texttt{complete\_\allowbreak{}UhatLt}).

\runinhead{Lower-bounded hatted Until.}
At every position where $F$ holds, the canonical policy resets the shared
clock (Section~\ref{subsec:canonical-main}).  The proof propagates an active obligation until a position satisfies
\[
  x\ge d
  \quad\text{and}\quad
  \overline{P}\Until \overline{Q},
\]
which yields the finite branch of $\Gamma$.  If no such position is reached,
the active-obligation invariant implies that $p$ remains continuously true,
and since every later position again carries a pending obligation that
requires a later $q$, the proposition $q$ recurs infinitely often.  This gives
\[
  \Always \overline{P}\land\Always\Eventually \overline{Q},
\]
the second branch of $\Gamma$.  This is formalized by
\texttt{complete\_\allowbreak{}UhatGe}, and by
\texttt{complete\_\allowbreak{}UhatGt} for the strict threshold $x>d$.

\runinhead{Upper-bounded hatted Release.}
The canonical policy resets the clock at every position where $F$ holds.  The
proof follows the currently active nontrivial Release obligation.  Until it
is discharged, the clock measures elapsed time from the relevant start and
$q$ continues to hold.  The weak until terminates either because $p\land q$
discharges the Release or because $x>d$, meaning that the upper-bounded window
has ended.  This is the role of \texttt{complete\_\allowbreak{}RhatLe}, and of
\texttt{complete\_\allowbreak{}RhatLt} for the strict window end $x\ge d$.

\runinhead{Lower-bounded hatted Release.}
Here the invariant is
\[
  \mathsf{RRes}
  \bigl(n,\CVal_{\rho^\star}(n,x),d,p,q\bigr).
\]
By Section~\ref{subsec:canonical-main}, the shared clock is preserved exactly
while $x<d$, the residual holds, and $p$ does not occur; in that case the
residual shifts
from $(i,v)$ to
\[
  (i+1,v+\Delta_i).
\]
The Coq lemma \texttt{rge\_\allowbreak{}residual\_\allowbreak{}shift} establishes this one-step update.
A converse lemma on the preservation policy recovers the residual information needed
at the next position, so the left side of the weak until remains valid until
either an ordinary untimed Release is established or $p$ occurs while the
clock is still below $d$.  This yields \texttt{complete\_\allowbreak{}RhatGe}.
The strict case uses $\mathsf{RRes}^{>}$, the pre-threshold test $x\le d$, and
the shift lemma \texttt{rgt\_\allowbreak{}residual\_\allowbreak{}shift}
(\texttt{complete\_\allowbreak{}RhatGt}).

As with soundness, ordinary timed operators need no additional completeness
lemmas: unfolding their definitions reduces them to Boolean structure and one
of the eight hatted primitive cases above.
